\section{Threat Model}
\label{sec:threat}

\subsection{Security Motivation}
Cache side-channel attacks, such as Prime+Probe and
Reload+ Refresh, rely on predicting or controlling cache
evictions. Understanding the replacement policy allows
attackers to construct reliable eviction sequences, while
incorrect assumptions can reduce attack effectiveness.
Our project investigates whether replacement behavior can
be learned without knowledge of the underlying hardware
implementation.

\subsection{Reverse-Engineering Setting}
We do not model an attacker. Instead, we reverse-engineer
the replacement policy of processors we can run code on,
with the access of an ordinary user. We know the cache
geometry but not its replacement policy. We can generate
memory-access sequences on our own memory and observe
eviction outcomes through cache hits and misses. Our goal
is to predict which cache line will be evicted based on
the access history and to identify the closest known policy.

\subsection{Simulation Assumptions}
Our simulator models a single $W$-way cache set with $W+1$
available memory blocks, or more when histories include
misses. Each experiment starts from a known
state, and eviction outcomes are observed without timing
noise, prefetching, or interference from other processes.
The simulation is therefore noise-free, and the full cache
state, including which block occupies each way and the
replacement bits or counters, is visible at every step.
Synthetic noise may be introduced later to evaluate model
robustness.

\subsection{Scope and Limitations}
Our simulation focuses on learning replacement behavior rather
than performing actual attacks. Eviction-set construction,
address mapping, cache-slice hashing, and hardware measurements
are excluded from this simulation stage but may be addressed
during later hardware experiments.


\section{Our Approach}
\label{sec:approach}

\subsection{Cache Simulation and Data Generation}

We train a GRU on sequences of memory accesses and their known eviction outcomes from cache simulations, allowing it to predict which cache line will be evicted for a new access sequence.

Our initial experiments focus on tree-based pseudo-LRU
(PLRU). We implemented a simulator for a single cache set
with associativities of 4, 8, and 16 ways. Each experiment
begins by filling the set with known blocks, followed by a
random sequence of accesses to those blocks. An additional block is then
accessed to force an eviction, which provides the training
label.

Each sample contains the access history, its length, and
the evicted way. History accesses are drawn uniformly from
the $W$ filled blocks, so every access is a hit and the
victim depends only on the policy's state. We generated
200,000 training samples for 4-way and 8-way caches and
1 million for 16-way caches, which have more state to learn.
Training histories contain between 0 and 32 accesses.

Samples are assigned to training, validation, and test sets
(90/5/5) by hashing the history, so no history appears in
more than one set. We also generated three held-out test
sets: histories of 33--128 and 129--512 accesses, to test
whether the model generalizes beyond the training lengths,
and 20,000 diagnostic histories where PLRU and LRU choose
different victims, to confirm the model has not simply
learned LRU. Figure~\ref{fig:sample} shows example inputs
and outputs.

\begin{figure}[htbp]
\centering
\footnotesize
\begin{tabular}{@{}l@{}}
\textbf{(a) Tree-PLRU} \\[2pt]
Input: $\underbrace{\texttt{FILL b0 b1 b2 b3}}_{\mbox{\scriptsize fill the set}}
\ \texttt{SEP}\
\underbrace{\texttt{b3 b2 b3 b0 b3 b1 b1}}_{\mbox{\scriptsize history (all hits)}}
\ \texttt{NEW}$ \\[14pt]
Output: way 2 is evicted, giving the training label \\[2pt]
\quad \begin{tabular}[t]{|c|c|c|c|}
\hline
way 0 & way 1 & way 2 & way 3 \\
\hline
0 & 0 & \textbf{1} & 0 \\
\hline
\end{tabular}
\quad (1 marks the evicted way) \\[8pt]
\textbf{(b) DRRIP} \\[2pt]
Input: $\texttt{FILL}\
\underbrace{\texttt{SRRIP}}_{\mbox{\scriptsize mode}}
\ \texttt{b0 b1 b2 b3 SEP}\
\underbrace{\texttt{b1 b4 b2}}_{\mbox{\scriptsize b4 misses}}
\ \texttt{NEW}$ \\[14pt]
Output: way 3 is evicted in 16 of 16 runs \\[4pt]
Input: \texttt{FILL BRRIP b0 b1 b2 b3 SEP b1 b4 b2 NEW} \\[2pt]
Output: way 0 (the new block b4) is evicted in most runs \\
\end{tabular}
\caption{Example model inputs and outputs for a 4-way set.
\texttt{FILL} precedes the blocks that fill the set,
\texttt{SEP} separates the fill from the access history,
and \texttt{NEW} marks the access of a new block, where the
model predicts which way is evicted. In (b), the same
history evicts different ways depending on the mode token.}
\label{fig:sample}
\end{figure}

\subsection{Neural Network Architecture and Training}

We use a Gated Recurrent Unit (GRU) network to learn
replacement behavior from memory-access sequences. Our
model consists of a 32-dimensional embedding layer, a
two-layer GRU with 128 hidden units, and a classification
layer that predicts the eviction victim.

We train the model by comparing its predictions to the correct eviction outcomes and adjusting it to improve accuracy. We gradually increase the length of the memory-access sequences during training and stop once the model consistently achieves 99.9\% accuracy on validation data. All initial experiments were conducted on a laptop CPU.

\subsection{Extension to an Adaptive Policy}

Intel LLCs are expected to use adaptive policies, so we
also simulated a 16-way DRRIP cache set. DRRIP switches
sets between two RRIP variants, SRRIP and BRRIP, which
differ only in how new blocks are inserted. BRRIP is
randomized, so each history was simulated 16 times and the
model was trained on the resulting victim distribution.
The global mode selection is not visible from one set, so
we gave the model the active mode as an extra input token
(Figure~\ref{fig:sample}b).
With hits-only histories both modes chose the same victim
97\% of the time, so these histories also access up to 8
additional blocks, causing misses. We generated 500,000
training histories, each labeled under both modes, and
trained for ten epochs. To compare policies on identical
inputs, we also labeled the same histories with tree-PLRU
and trained the same model on them.

\subsection{Hardware Characterization}

After validating our approach in simulation, we plan to
collect eviction measurements from a local Intel Celeron
N3450 (Goldmont). It has no L3, so its L2 is the
last-level cache.
For a $W$-way cache set, we will select $W+1$ memory blocks
mapping to the same set. We will fill the cache, execute
a generated access history, and introduce an additional
block to force an eviction. These $W+1$ blocks suffice for
hits-only histories. Adaptive and counter-based policies
require histories with misses, so for these the history
will draw from up to $W+8$ same-set blocks, as in our DRRIP
simulation, plus one new block for the final access.

To identify the evicted block, we will replay the experiment
and probe each block individually using cache hit/miss
measurements. Replaying is necessary because each probe can
modify the cache state. With hits-only histories, only the
original $W$ blocks need to be probed. When histories include
misses, blocks can also be evicted during the history, so we
will probe every block both with and without the final
access; the victim is the block that is resident before the
final access but not after it.

On hardware we observe which block was evicted, not which
way it occupied. With hits-only histories the two are the
same, because block $b_i$ stays in way $i$. With misses,
blocks move between ways, so for these histories our labels
will be the evicted block. We will start with hits-only
histories, which are cheaper to measure, and add histories
with a few misses only where candidate policies cannot
otherwise be distinguished.

Once validated on the Celeron, we plan to apply
the same methodology to the L2 and, if time permits, the
Last-Level Cache (LLC) of an LRC machine.

\subsection{Replacement Policy Identification}

We will fine-tune the GRU on hardware measurements so that
it predicts the hardware's evictions, and check its accuracy
on held-out histories and cache sets against how often
repeated measurements agree. We will then compare the
model's predicted victims against those of each known
replacement-policy simulator. A policy matches if it agrees
on at least 95\% of diagnostic histories and clearly more
than every other candidate. We also generated diagnostic access
histories where candidate policies predict different
victims, allowing us to distinguish between policies
that otherwise behave similarly.

The trained model can be queried on any number of new
histories at no hardware cost, so we can search for
histories that separate the leading candidates and verify
only those on hardware. It also smooths over measurement
noise, and if no known policy matches, it still provides a
working model of the unknown policy to analyze. We will
also compare hardware measurements directly against each
simulator, which does not depend on model accuracy.